\chapter{Introduction}

\section{Overview}

Vehicle-to-Vehicle (V2V) communication is a key enabler for cooperative safety functions such as collision
warnings, hazard propagation, and context-aware driver assistance. DSRC (IEEE 802.11p) and cellular V2X
(C-V2X) remain the primary technologies for deployment and standardization \cite{secure_iov,lorawan_v2roadside}.
This project investigates LoRa as a low-cost experimental platform for decentralized dissemination in
infrastructure-constrained settings, not as a replacement for purpose-built vehicular standards.

The proposed LoRa-based V2V mesh system uses a three-layer stack, with detailed responsibilities and design
rationale presented in Chapter 3.

This report documents the complete final-year project workflow from design rationale to implementation and
validation. Results are presented through complementary simulation and real-node experiments, followed by a
consolidated discussion of implications and limitations.

\section{Motivation}

The core question is whether a \textit{fully decentralized} vehicular mesh can provide useful multi-hop
safety-message dissemination using only peer forwarding and local decision-making. This focus is relevant to
academic prototyping and to scenarios where infrastructure support may be sparse, unavailable, or temporarily
disrupted \cite{lora_mesh_emergency,lora_mesh_iot}.

\section{Problem Definition}

This work evaluates whether a \textit{fully decentralized} LoRa-based V2V mesh can provide reliable and
efficient multi-hop safety-message dissemination for infrastructure-free operation.

The problem is challenging because LoRa links impose strict communication constraints, including low data rate,
long Time-on-Air, duty-cycle limits, and shared-medium contention \cite{lora_mesh_iot,routing_mesh_2020}.
Without careful protocol control, these constraints can lead to excessive collisions, redundant rebroadcast,
and unstable dissemination behavior in dynamic vehicular topologies \cite{vdtn_routing_eval,lora_mesh_emergency}.
Hence, the central problem is to design a decentralized protocol stack that preserves dissemination reach while
keeping channel usage bounded and operational behavior predictable.

\section{Objectives}

The project objectives are:

\begin{itemize}
	\item Implement a three-layer architecture with clear separation between application logic, forwarding control,
	and radio-facing transmission services.
	\item Implement decentralized network-layer dissemination using managed flooding, duplicate suppression,
	hop-bounded propagation, and policy-aware relay timing.
	\item Implement link-layer channel-aware behavior, including sensing-driven transmission admission and
	optional reliability signaling for critical traffic classes.
	\item Demonstrate that the system can deliver safety messages across multiple hops with bounded redundancy,
	without dependence on RSUs, fixed gateways, or centralized coordination.
\end{itemize}

\section{Contributions}

Experimental results confirm end-to-end message delivery across nodes beyond direct transmission range, correct
priority ordering with an observed delay ratio consistent with the intended fourfold scaling, a 33\% reduction
in relay events under directional mode, and distributed best-relay selection without explicit inter-node
negotiation.

Simulation and real-node evaluations provided consistent evidence of intended system behavior. The simulation
campaign verified forwarding correctness, bounded propagation policies, and timing consistency under modeled
channel impairments. Four-node field testing demonstrated multi-hop reach extension, priority-aware relay
ordering, directional propagation control, and autonomous suppression of redundant relays.

